\documentclass{Phyport}

\usepackage{hyperref}
\hypersetup{
    colorlinks=true,
}
% 选用宏（如有额外需要的宏包请在此添加）
\usepackage{float}       % 提供 [H] 选项
\usepackage{booktabs}    % 三线表
\usepackage{amsmath}     % 数学公式
\usepackage{caption}     % 标题设置
\usepackage{siunitx}     % SI 单位排版
 %选用宏

\usepackage{circuitikz}
% 双臂电桥（开尔文电桥）原理电路图
% 拓扑：R1(Px1→B)、R2(B→Ps2) 与主回路 Rx-r-Rs 并联；
%       R3(Px2→D)、R4(D→Ps1) 与粗导线 r 并联；检流计 G 接在 B、D 之间。
%       平衡时由 R1 I1 = Rx I3 + R3 I2、R2 I1 = Rs I3 + R4 I2 得 Rx = (R1/R2) Rs。
\newcommand{\ArmBridge}{%
\begin{circuitikz}[european,line width=0.7pt]

  % ---------- 顶部主回路：四端电阻 Rx —— 粗导线 r —— 比较臂 Rs ----------
  \draw (1.5,4.8) rectangle (5.1,5.2);
  \node[font=\small] at (3.3,5.0) {$R_x$};
  \draw (6.4,4.8) rectangle (10.0,5.2);
  \node[font=\small] at (8.2,5.0) {$R_s$};
  \draw[line width=2pt] (5.1,5.0) -- (6.4,5.0);
  \node[below,font=\footnotesize] at (5.75,4.94) {$r$};

  % ---------- 比率臂：R1/R2、R3/R4 各留出间隙 ----------
  \draw (2.1,4.8) -- (2.1,4.4);
  \draw (2.1,4.4) to[R] (2.1,3.4);
  \draw (2.1,3.4) -- (2.1,2.8);
  \draw (2.1,2.8) to[R] (2.1,1.8);
  \draw (4.5,4.8) -- (4.5,4.4);
  \draw (4.5,4.4) to[R] (4.5,3.4);
  \draw (4.5,3.4) -- (4.5,2.8);
  \draw (4.5,2.8) to[R] (4.5,1.8);

  % ---------- 检流计 ----------
  \node[circle,draw,inner sep=1pt,font=\small] (G) at (3.3,3.1) {$G$};
  \draw (2.1,3.1) -- (G.west);
  \draw (G.east) -- (4.5,3.1);
  \fill (2.1,3.1) circle (1pt);
  \fill (4.5,3.1) circle (1pt);

  % ---------- 电位端回路 ----------
  \draw (2.1,1.8) -- (2.1,0.6) -- (9.4,0.6) -- (9.4,4.8);   % R2 → Ps2
  \draw (4.5,1.8) -- (4.5,1.2) -- (7.0,1.2) -- (7.0,4.8);   % R4 → Ps1

  % ---------- 直流电源 ----------
  \draw (1.5,5.0) -- (-0.4,5.0) -- (-0.4,-1.3) -- (4.0,-1.3);
  \draw (4.0,-1.52) -- (4.0,-1.08);
  \draw (4.45,-1.42) -- (4.45,-1.18);
  \draw (4.45,-1.3) -- (10.8,-1.3) -- (10.8,5.0) -- (10.0,5.0);
  \node[below,font=\small] at (4.22,-1.72) {$\varepsilon$};

  % ---------- 端子与元件标注 ----------
  \node[above left,font=\footnotesize]  at (1.45,5.06)  {$C_{x1}$};
  \node[above,font=\footnotesize]       at (4.95,5.22)  {$C_{x2}$};
  \node[above,font=\footnotesize]       at (6.55,5.22)  {$C_{s1}$};
  \node[above right,font=\footnotesize] at (10.05,5.06) {$C_{s2}$};
  \node[right,font=\footnotesize]       at (2.16,4.6)   {$P_{x1}$};
  \node[right,font=\footnotesize]       at (4.56,4.6)   {$P_{x2}$};
  \node[right,font=\footnotesize]       at (7.06,4.6)   {$P_{s1}$};
  \node[right,font=\footnotesize]       at (9.46,4.6)   {$P_{s2}$};
  \node[left,font=\footnotesize]        at (1.86,3.9)   {$R_1$};
  \node[left,font=\footnotesize]        at (1.86,2.3)   {$R_2$};
  \node[right,font=\footnotesize]       at (4.74,3.9)   {$R_3$};
  \node[right,font=\footnotesize]       at (4.74,2.3)   {$R_4$};
  \node[left,font=\footnotesize]        at (2.0,3.1)    {$B$};
  \node[right,font=\footnotesize]       at (4.6,3.1)    {$D$};

\end{circuitikz}}


\exname{直流双臂电桥} %实验名称
\extable{5} %实验桌号
\instructor{宋斌} %指导教师
\class{电科2502} %班级
\name{管云哲} %姓名
\stuid{3250104519} %学号

\nyear{2026} %年
\nmonth{10} %月
\nday{13} %日
\nweekday{二} %星期几

\redate{} %如有实验补做，补做日期
\resitu{} %情况说明：

\begin{document}
\setcounter{page}{1}
\makecover

\section{预习报告（10分）}
（注：将已经写好的“物理实验预习报告”内容拷贝过来）

\subsection{实验综述（5分）}
（自述实验现象、实验原理和实验方法，包括必要的光路图、电路图、公式等。不超过500字。）

实验现象：电桥平衡时，检流计指针指零。

实验原理：用单臂电桥测低电阻时，接线电阻与接触电阻（$0.01\sim1\,\Omega$）会与待测电阻 $R_x$ 串联而被计入测量结果，误差极大。为此将低电阻改为四端接入：电流端 $C_{x1}$、$C_{x2}$ 通电流，电位端 $P_{x1}$、$P_{x2}$ 取电压，只有两电位端之间的电阻才是实测值，附加电阻被“转移”至待测电阻之外。把两个四端电阻 $R_x$、$R_s$ 接入桥路，并另设阻值均大于 $10\,\Omega$ 的标准电阻 $R_3$、$R_4$，用 $r<0.001\,\Omega$ 的粗导线连接 $R_x$ 与 $R_s$，即构成双臂电桥，其电路如图所示。

\begin{figure}[H]
    \centering
    \ArmBridge
    \caption{双臂电桥（开尔文电桥）原理电路图}
    \label{fig:armbridge}
\end{figure}

平衡时 $I_G=0$，由基尔霍夫定律可得
\begin{equation}
    R_x=\frac{R_1}{R_2}R_s+\frac{R_4 r}{R_3+R_4+r}\left(\frac{R_1}{R_2}-\frac{R_3}{R_4}\right)
    \label{eq:rx}
\end{equation}
若使两对比率臂始终保持 $R_1/R_2=R_3/R_4$，则第二项为零，待测低电阻为
\begin{equation}
    R_x=\frac{R_1}{R_2}R_s
    \label{eq:rx2}
\end{equation}
金属的电阻随温度线性变化：$R_t=R_0(1+\alpha t)$，$R_0$ 为 $0\,^\circ\mathrm{C}$ 时的阻值，$\alpha$ 为电阻温度系数。

实验方法：（1）测电阻率。接通检流计并调零，灵敏度置于最低；接入待测金属导体，选适当倍率，先按 B 通电、后按 G 检流，调节粗调与细调盘使检流计偏转减小，逐步提高灵敏度至最高档调平衡，读数 $R_x=$ 倍率 $\times$（粗调$+$细调）；再用游标卡尺测导体直径 $d$、用米尺测长度 $L$。（2）测电阻温度系数。在最高灵敏度下调平衡，开启加热装置设定温度，用升温或降温法每改变约 $5\,^\circ\mathrm{C}$ 测一次 $R$ 与 $t$，记录 10 组数据。

\subsection{实验重点（3分）}
（简述本实验的学习重点，不超过100字。）

掌握双臂电桥测低电阻的原理，理解四端接法与双比率臂使 $R_1/R_2=R_3/R_4$ 以消除引线电阻、接触电阻影响的方法；学会 QJ44 型直流双臂电桥的操作，并由 $R=\dfrac{4L}{\pi d^2}\rho$ 测定金属电阻率与电阻温度系数。

\subsection{实验难点（2分）}
（简述本实验的实现难点，不超过100字。）

操作顺序须“先按 B 后按 G”，并随平衡程度逐步提高灵敏度，否则检流计易过载；高灵敏度下平衡点难以判准，需兼顾调节速率与读数精度；升温测量时导体温度与显示温度不一致，须待温度稳定后同步读取温度与电阻。

\section{实验操作与原始数据（50分）}
（将有老师签名的“自备数据记录草稿纸”的扫描或手机拍摄图粘贴在下方，完整保留姓名，学号，教师签字和日期。按老师的要求拍摄实验过程图，附在此处）

\subsection{实验仪器}

QJ44 型直流双臂电桥（0.2 级，含检流计与倍率、粗调、细调盘），温控加热装置，待测铜导体棒，游标卡尺（分度值 $0.02\,\mathrm{mm}$）。

\subsection{实验步骤}

\begin{enumerate}
    \item 接通电桥电源与检流计，将灵敏度调至最低并调零检流计。
    \item 按四端接法接入待测铜导体棒，选择适当倍率；先按 B 通电，后按 G 检流，调节粗调盘与细调盘使检流计偏转逐渐减小。
    \item 逐步提高电桥灵敏度，反复微调至检流计指零，在最高灵敏度下读数 $R_x=$ 倍率 $\times$（阻值粗调 $+$ 阻值细调）。
    \item 用游标卡尺在导体棒不同位置测量直径 $d$，用米尺测量两电位端之间的长度 $L$。
    \item 保持电桥平衡状态接入加热装置，设定目标温度，待温度稳定后每约 $5\,^\circ\mathrm{C}$ 同步记录温度 $t$ 与电阻 $R$，共 10 组。
\end{enumerate}

\begin{figure}[H]
    \centering
    \includegraphics[width=0.95\linewidth]{实验过程}
    \caption{实验装置与接线（SR-1/a 型双臂电桥四端接线夹具、QJ44 型直流双臂电桥、温控加热装置）}
    \label{fig:process}
\end{figure}

\begin{figure}[H]
    \centering
    \includegraphics[width=0.95\linewidth]{实验数据}
    \caption{原始数据记录（含姓名、学号、教师签字与日期）}
    \label{fig:rawdata}
\end{figure}

\section{结果与分析（30分）}
\subsection{数据处理与结果（20分）}
（列出数据表格、选择适合的数据处理方法、写出测量或计算结果。）

\subsubsection{测量金属导体的电阻率}

\begin{table}[H]
    \centering
    \caption{金属导体电阻率测量数据}
    \label{tab:resistivity}
    \begin{tabular}{cccc}
        \toprule
        阻值 $R$ / $\Omega$ & 直径 $d$ / mm & 长度 $L$ / cm & 电阻率 $\rho$ / $\Omega\cdot\mathrm{m}$ \\
        \midrule
        $6.52\times10^{-4}$ & 4.04 & 30.00 & $2.79\times10^{-8}$ \\
        \bottomrule
    \end{tabular}
\end{table}

电阻率由 $R=\dfrac{4L}{\pi d^2}\rho$ 反解得到：
\begin{equation}
    \rho=R\frac{\pi d^2}{4L}
    =6.52\times10^{-4}\times\frac{\pi\times(4.04\times10^{-3})^2}{4\times30.00\times10^{-2}}
    =2.786\times10^{-8}\ \Omega\cdot\mathrm{m}
    \label{eq:rho}
\end{equation}

\subsubsection{测量金属导体的电阻温度系数}

\begin{table}[H]
    \centering
    \caption{电阻随温度的变化}
    \label{tab:temperature}
    \begin{tabular}{ccccccccccc}
        \toprule
        次数 & 1 & 2 & 3 & 4 & 5 & 6 & 7 & 8 & 9 & 10 \\
        \midrule
        $t$ / $^\circ\mathrm{C}$ & 24.3 & 29.3 & 34.3 & 39.3 & 44.3 & 49.3 & 54.3 & 59.3 & 64.3 & 69.3 \\
        $R_x$ / $10^{-3}\Omega$ & 4.815 & 4.908 & 5.007 & 5.102 & 5.190 & 5.288 & 5.359 & 5.471 & 5.566 & 5.660 \\
        \bottomrule
    \end{tabular}
\end{table}

\textbf{处理方式 1：分组逐差法。} 将数据按 $i$ 与 $i+5$ 分两组代入
\begin{equation}
    \alpha_i=\frac{R_{x(i+5)}-R_{x(i)}}{R_{x(i)}t_{(i+5)}-R_{x(i+5)}t_{(i)}},\qquad i=1,2,3,4,5
    \label{eq:alphai}
\end{equation}
得
\begin{equation*}
    \alpha_1=4.34\times10^{-3}\,^\circ\mathrm{C}^{-1},\quad
    \alpha_2=4.12\times10^{-3}\,^\circ\mathrm{C}^{-1},\quad
    \alpha_3=4.25\times10^{-3}\,^\circ\mathrm{C}^{-1},
\end{equation*}
\begin{equation*}
    \alpha_4=4.24\times10^{-3}\,^\circ\mathrm{C}^{-1},\quad
    \alpha_5=4.31\times10^{-3}\,^\circ\mathrm{C}^{-1}
\end{equation*}
取平均得 $\bar{\alpha}=\dfrac{1}{5}\sum_{i=1}^{5}\alpha_i=4.25\times10^{-3}\,^\circ\mathrm{C}^{-1}$。

\textbf{处理方式 2：最小二乘线性拟合（图解法）。} 作 $R_t-t$ 曲线如图~\ref{fig:rtcurve}，线性拟合得
\begin{equation*}
    R_t=(0.01867\,t+4.36300)\times10^{-3}\ \Omega
\end{equation*}

\begin{figure}[H]
    \centering
    \includegraphics[width=0.78\linewidth]{Rt_curve}
    \caption{铜导体电阻 $R_t$ 随温度 $t$ 的变化及线性拟合}
    \label{fig:rtcurve}
\end{figure}

由 $R_t=R_0(1+\alpha t)$ 知斜率与截距之比即为温度系数：
\begin{equation*}
    \alpha_0=\frac{0.01867}{4.36300}\ ^\circ\mathrm{C}^{-1}=4.28\times10^{-3}\ ^\circ\mathrm{C}^{-1}
\end{equation*}
拟合直线的截距 $R_0=4.363\times10^{-3}\,\Omega$ 即导体在 $0\,^\circ\mathrm{C}$ 时的电阻。两种处理方法所得温度系数一致（$4.25\times10^{-3}$ 与 $4.28\times10^{-3}\,^\circ\mathrm{C}^{-1}$）。

\subsection{误差分析（10分）}
（运用测量误差、相对误差或不确定度等分析实验结果，写出完整的结果表达式，并分析误差原因。）

\subsubsection{测量金属导体的电阻率}

由式~\eqref{eq:rho}，$\rho=R\dfrac{\pi d^2}{4L}$，先求各直接测量量的不确定度。电桥为 0.2 级，按均匀分布取
\begin{equation*}
    u(R)=\frac{0.2\%\times R}{\sqrt{3}}=\frac{0.2\%\times6.52\times10^{-4}}{\sqrt{3}}=7.5\times10^{-7}\ \Omega
\end{equation*}
游标卡尺与米尺的分度值分别为 $0.02\,\mathrm{mm}$、$0.5\,\mathrm{mm}$，同样按均匀分布：
\begin{equation*}
    u(d)=\frac{0.02}{\sqrt{3}}\,\mathrm{mm}=0.012\,\mathrm{mm},\qquad
    u(L)=\frac{0.5}{\sqrt{3}}\,\mathrm{mm}=0.29\,\mathrm{mm}
\end{equation*}
作间接测量不确定度合成，电阻率的相对不确定度为
\begin{equation}
    \frac{u(\rho)}{\rho}=\sqrt{\left(\frac{u(R)}{R}\right)^2+\left(2\frac{u(d)}{d}\right)^2+\left(\frac{u(L)}{L}\right)^2}
    =\sqrt{(0.115\%)^2+(0.572\%)^2+(0.096\%)^2}=0.59\%
    \label{eq:relrho}
\end{equation}
故
\begin{equation*}
    u(\rho)=\rho\cdot\frac{u(\rho)}{\rho}=2.786\times10^{-8}\times0.59\%=0.017\times10^{-8}\ \Omega\cdot\mathrm{m}
\end{equation*}
电阻率的测量结果为
\begin{equation*}
    \rho=(2.79\pm0.02)\times10^{-8}\ \Omega\cdot\mathrm{m}
\end{equation*}

由式~\eqref{eq:relrho} 可知，三项相对分量中直径项 $2u(d)/d=0.572\%$ 占绝对主导，长度项与电桥项分别仅 $0.096\%$ 与 $0.115\%$，故提高测量精度的关键在于直径的测量。误差主要来源：

\begin{itemize}
    \item 导体棒的电流接头可在小范围内晃动，导致两电位端之间的有效长度 $L$ 难以准确界定；
    \item 游标卡尺读数存在估读误差，且导体棒存在微小弯曲，其截面并非理想圆，同一位置不同方向的直径读数有差异；
    \item 电桥本身的 0.2 级准确度限制，以及接触电阻未被比率臂完全抵消的残余影响。
\end{itemize}

\subsubsection{测量金属导体的电阻温度系数}

两种处理方法所得 $\alpha$ 分别为 $4.25\times10^{-3}$ 与 $4.28\times10^{-3}\,^\circ\mathrm{C}^{-1}$，与理论值 $\alpha=4.33\times10^{-3}\,^\circ\mathrm{C}^{-1}$ 相比，相对误差分别为
\begin{equation*}
    E_1=\frac{|\bar{\alpha}-\alpha|}{\alpha}=1.8\%,\qquad
    E_2=\frac{|\alpha_0-\alpha|}{\alpha}=1.2\%
\end{equation*}
两者均在 2\% 以内相符。误差主要来源：

\begin{itemize}
    \item 加热过程中导体自身温度分布不均，其真实温度与温控装置显示的温度并不完全一致；
    \item 温度稳定后调节电桥平衡需要时间，在此期间导体温度会有微小漂移，导致 $R$ 与 $t$ 不能严格同步；
    \item 拟合直线的截距 $R_0$ 是由 $24\sim69\,^\circ\mathrm{C}$ 数据外推至 $0\,^\circ\mathrm{C}$ 得到的，外推距离远，截距的相对误差被放大后传递到 $\alpha$。
\end{itemize}

\section{实验探讨与思考题（10分）}
（对实验内容、现象和过程的小结，不超过100字。并解答教材或讲义或老师布置的思考题，请先写题干，再作答。）

本次实验用 QJ44 型双臂电桥测铜导体的低电阻，得到电阻率 $\rho=(2.79\pm0.02)\times10^{-8}\,\Omega\cdot\mathrm{m}$，温度系数 $\alpha\approx4.3\times10^{-3}\,^\circ\mathrm{C}^{-1}$，与理论值相符。实验中体会到四端接法与双比率臂 $R_1/R_2=R_3/R_4$ 对消除引线及接触电阻的关键作用，也理解了双臂电桥相对惠斯登电桥在低电阻测量上的优势。

\subsection*{思考题 1：双臂电桥和惠斯登电桥有哪些异同？}

相同点：两者都利用电桥平衡时检流计指零、由桥臂电阻比与标准电阻反推待测电阻，都是比较法测电阻；均由改变桥臂电阻来调节平衡。

不同点：惠斯登电桥只有一个桥臂接入待测电阻，接线电阻和接触电阻与 $R_x$ 串联后直接进入结果，测 $1\,\Omega$ 以下的低电阻时误差极大，因此只适用于 $1\,\Omega\sim1\,\mathrm{M}\Omega$ 的中值电阻；双臂电桥有两个对比臂同时接入待测电阻 $R_x$ 与比较臂 $R_s$，并采用四端接法与双比率臂，能把引线电阻和接触电阻“赶出”测量回路，故用于测量 $1\,\Omega$ 以下的低电阻。

\subsection*{思考题 2：为什么双臂电桥测量低电阻时能够消除（或减小）附加电阻对测量结果的影响？}

由式~\eqref{eq:rx}，$R_x$ 中包含附加项 $\dfrac{R_4r}{R_3+R_4+r}\left(\dfrac{R_1}{R_2}-\dfrac{R_3}{R_4}\right)$，减小附加电阻的影响即要使其尽可能小，这由两点保证：

\begin{enumerate}
    \item 连接 $R_x$ 与 $R_s$ 的粗导线电阻 $r<0.001\,\Omega$，而 $R_3$、$R_4$ 均大于 $10\,\Omega$，故因子 $\dfrac{R_4r}{R_3+R_4+r}$ 极小；
    \item 调节时始终保持 $\dfrac{R_1}{R_2}=\dfrac{R_3}{R_4}$，使括号内的差为零。此时 $R_x=\dfrac{R_1}{R_2}R_s$，与 $r$ 无关。
\end{enumerate}

此外，电流端 $C_{x1}$、$C_{x2}$、$C_{s1}$、$C_{s2}$ 处的接触电阻只与电源回路串联，电位端 $P_{x1}$、$P_{x2}$、$P_{s1}$、$P_{s2}$ 处的接触电阻只与高阻值的比率臂串联，二者均不改变平衡条件，从而消除了附加电阻对测量结果的影响。

\subsection*{思考题 3：如果四端电阻的电流端和电位端接反了，对测量有哪些影响？}

接反后电流端与被测段串联的附加电阻（接触电阻与引线电阻）不再被排除在测量回路之外，而 $R_3$、$R_4$ 支路则串入了本应属于电流回路的电阻，式~\eqref{eq:rx} 中括号内的差不再能为零，附加电阻将直接进入 $R_x$ 的测量结果，使测得的电阻显著偏大；对于阻值本身只有 $10^{-4}\,\Omega$ 量级的低电阻，这一偏差甚至可能超过待测电阻本身，测量失去意义。

{\fangsong 注意事项：
\begin{enumerate}[label=\arabic*.]
    \item 用PDF格式上传“实验报告”，文件名：学生姓名+学号+实验名称+周次。
    \item “实验报告”必须递交在“学在浙大”本课程内对应实验项目的“作业”模块内。
    \item “实验报告”成绩必须在“浙江大学物理实验教学中心网站”-“选课系统”内查询。
    \item 教学评价必须在“浙江大学物理实验教学中心网站”-“选课系统”内进行，学生必须进行教学评价，才能看到实验报告成绩，教学评价须在本次实验结束后3天内进行。
\end{enumerate}}
\begin{center}
    \fbfangsong{浙江大学物理实验教学中心制}
\end{center}
\end{document}
